\chapter{Weierstrass preparation}

%% We are going to reformat some of this file by generalising the topological nilpotence and
%% Power bounded defs to as in Wedhorn (see PowerBounded file)
%% The stuff on units etc can then be proven in generality (I think)

When computing zeroes of a power series it would be convenient if we could instead reduce it to
calculating zeroes of a polynomial. This turns out to be possible for restricted power series using
the $p$-adic Weierstrass preparation. This section is dedicated to proving this classical result. We
will follow \ref{BGR}.

\newcommand{\pb}[1]{\overset{\circ}{#1}}
\newcommand{\tn}[1]{\check{#1}}
\newcommand{\res}[1]{\widetilde{#1}}
\newcommand{\quot}[2]{{\raisebox{.2em}{$#1$}\left/\raisebox{-.2em}{$#2$}\right.}}

\section{Notation}

Before we state the theorem we first need to set some notation. Let $R$ be a semi-normed ring.

\begin{definition}
  We say an element $r$ is \emph{topologically nilpotent} if
    \[
     \lim r^n = 0.
    \]
    Denote the set of all topologically nilpotent elements of $R$ as $\tn{R}$.
\end{definition}

Unfortunately, the set $\tn{R}$ is not, in general, an ideal of $R$. However, we can define a
subring of $R$, that has it an ideal.

\begin{definition}
  We say an element $r$ is \emph{power-bounded} if the set
  \[
  \left\{ \lvert r^n \rvert : n \in \N \right\}
  \]
  is bounded.
  Denote the set of all power-bounded elements of $R$ as $\pb{R}$.
\end{definition}

\begin{lemma}
  $\pb{R}$ is a clopen subring of $R$, and $\tn{R}$ is an ideal in $\pb{R}$.
\end{lemma}

\begin{proof}
  Let $a, b \in \pb{R}$. Let $M$ be an upper bound for $\lvert a^n \rvert$ and $\lvert b^n \rvert$
  for all $n$. Then
  \[
  \lvert (a b )^n \rvert \leq \lvert a^n \rvert \lvert b^n \rvert \leq M^2,
  \]
  similarly we can use triangle inequality for closure of addition.
  For negation since, $R$ is a seminormed ring
  \[
  \lvert (- a)^n \rvert = \lvert a^n \rvert \leq M
  \]
  and thus since $1 \in \pb{R}$, we see it is a subring.

  If $a \in \nt{R}$ and $b \in \pb{R}$, then
  \[
  \lvert (a b)^n \rvert \leq \lvert a^n \rvert \lvert b^n \rvert \leq \lvert a^n \rvert M \to 0
  \]
  that is $ a b \in \tn{R}$, and we have shown $\tn{R}$ is an ideal.

  Finally, since $\pb{R}$ is a subgroup and contains $\tn{R}$, which is open in $R$, we have $\pb{R}$
  is clopen in $R$.
\end{proof}

Given these we can then define their quotient, $\res{R} = \pb{R} / \tn{R}$.

\section{Weierstrass division}

Weierstrass preparation will follow as an application of Weierstrass division, so we first prove that.
To do this we need to introduce the following terminology.


\begin{definition}
  Let $s$ be a natural number, we say a strictly convergent power series
  $g \in R \langle x_1 \dots, x_n \rangle >$ is $x_n$-distinguished
  of degree $s$ if, when writing $g = \sum_{v = 0}^\infty g_v(x_1, \dots, x_{n-1}) x_n^v$:
  \begin{itemize}
    \item $g_s$ is a unit; and
    \item $ \lvert g_s \rvert = \lvert g \rvert$ and $\lvert g_s \rvert > \lvert g_v \rvert \; \forall
      v > s$.
  \end{itemize}
\end{definition}

That is, $g_s$ achieves the gauss norm of $g$ and is the last index to do so.

\begin{theorem}
  Let $s$ be a natural number and $g \in T_n$ be $x_n$-distinguished of degree $s$. Then for each
  $f \in T_n$, there exists a unique $ q \in T_n$ and a unique $r \in T_{n-1}[x_n]$ with degree
  $r < s$ such that
  \[
    f = q g + r.
  \]
  Moreover, we have the following estimate
  \[
    \lvert f \rvert = \max \left\{ \lvert q \rvert \lvert g \rvert, \lvert r \rvert \right\},
  \]
  that is
  \[
    \lvert q \rvert \leq \lvert g \rvert^{-1} \quad \text{and} \quad \lvert r \rvert \leq \lvert f \rvert.
  \]
\end{theorem}

We break the proof up into the following steps.

\begin{definition}
  We say a subgroup $H$ of $G$ is $\epsilon$-dense if there is a real number $\epsilon < 1$ such that
  for each $g \in G$ there exists an $h \in H$ with
  \[
    \lvert g + h \rvert \leq \epsilon \lvert g \rvert.
  \]
\end{definition}

\begin{lemma}
  Let $(G, \lvert \cdot \rvert)$ be a normed group and $H$ a subgroup of $G$ $\epsilon$-dense in $G$.
  Then $H$ is dense in $G$.
\end{lemma}

To prove this we will use the following definition.

\begin{definition}
  Let $G$ be a seminormed group and let $H$ be any subgroup of $G$. For each $g \in G$, the nonnegative
  real number
  \[
  \lvert g, H \rvert := \inf_{h \in H} \lvert g + h \rvert
  \]
  is called the distance form $g$ to $H$.
\end{definition}

Using this definition it is clear to see the following:
\begin{itemize}
  \item $\lvert a + h, H \rvert = \lvert a, H \rvert$ for all $h \in H$;
  \item $\lvert g, H \rvert = \inf_{h \in H} d(a,h)$, where $d$ is the normal distance in a seminormed
    group; and
  \item $\overline{H} = \left\{ g \in G : \lvert g, H \rvert = 0 \right\}$.
\end{itemize}

\begin{proof}
  By the final bullet point above it suffices to show that $\lvert g, H \rvert = 0$ for all $g \in G$
  - as then $\overline{H} = G$, and we see $H$ is dense.

  Therefore, assume there is $g \in G$ such that $\lvert g, H \rvert > 0.$ Choose $h_1 \in H$ such that
  \[
    \lvert g + h_1 \rvert < \epsilon^{-1} \lvert g, H \rvert,
  \]
  this is possible since $\epsilon^{-1} > 1.$
  By $\epsilon$-dense, there exists $h_2$ such that
  \[
  \lvert (g + h_1) + h_2 \rvert \leq \epsilon \lvert g + h_1 \rvert.
  \]
  However, $h_1 + h_2 \in H$, thus the inequality obtained by combining the previous two,
  \[
    \lvert g + (h_1 + h_2) \rvert < \lvert g, H \rvert,
  \]
  is impossible. Therefore, $\lvert g, H \rvert = 0$ for all $g \in G$.
\end{proof}

With this we can then start proving parts of Weierstrass division.

\begin{lemma}\label{WD_bounds}
  Assume we are in the case of Weierstrass division and furthermore $\lvert g \rvert = 1$. Then, if
  there exists a representation
  \[
    f = q g + r
  \]
  we have the estimates
  \[
    \lvert q \rvert \leq \lvert f \rvert \quad \text{and} \quad \lvert r \rvert \leq \lvert f \rvert.
  \]
\end{lemma}

\begin{proof}
  Multiplying $f = q g + r$ by a scalar, we can assume
  \[
    \max \left\{ \lvert q \rvert, \lvert r \rvert \right\} = 1
  \]
  which implies $\lvert f \rvert \leq 1$ by the nonarchimedean property and $f = q g + r$.
  Thus, it suffices to show that $\lvert f = 1 \rvert.$ Assume the contrary, that is $\lvert f \rvert < 1$
  and so reduction into the residue field gives
  \[
    0 = \tilde{f} = \tilde{q} \tilde{g} + \tilde{r}.
  \]
  Since
  \[
    \deg \tilde{g} = s > \deg r > \deg \tilde{r}
  \]
  this would imply $\tilde{q} = \tilde{r} = 0$, which is a contradiction to the assumption that
  $\max {\lvert q \rvert, \lvert r \rvert} = 1$.
\end{proof}

\begin{lemma}\label{WD_uniqueness}
  Suppose we have a valid representation for Weierstrass division, then it is unique.
\end{lemma}

\begin{proof}
  Consider two representations
  \[
    f = q g + r \quad \text{and} \quad f = q' g + r',
  \]
  then we have
  \[
    0 = (q - q') g + (r - r').
  \]
  Applying the bounds to this new equation gives:
  \[
    \lvert q - q' \rvert \leq \lvert 0 \rvert = 0
  \]
  and similar for $r - r'$. Thus, $q - q' = 0 = r - r'$ and we see the two representations are equal.
\end{proof}


\begin{lemma}\label{WD_existence}
  There exists a representation
  \[
    f = q g + r
  \]
  in Weierstrass division.
\end{lemma}


%% I think to make sense of the epimorphism I will need prop 2 page 193
%% that proves \res{T}_n = \res{R}[n]... and specifically
%% \pb{T}_n = T_n^{\circ} ( = valuation ring )
%% \np{T}_n = T_n^\vee
%% \res{T}_n = T_n^\tilde

\begin{proof}
  Assume $g = \lvert 1 \rvert$ and consider the set
  \[
    B = \left\{ qg + r : r \in T_{n-1}[x_n], \; \deg r > s, \; q \in T_n \right\},
  \]
  this is a closed subgroup of $T_n$ and it suffices to show $B = T_n$.
  Writing $g = \sum g_s (x_1, \cdots, x_{n-1}) x_n^s$, define
  \[
    \varepsilon = \max_{s < t} \lvert g_t \rvert,
  \]
  since $g$ is $x_n$-distinguished of degree $s$ we must have $\varepsilon < 1$.

  Consider the natural ring epimorphism
  \[
    \tau_\varepsilon : \pb{T}_n \to \res{R}_\varepsilon [x_1, \cdots, x_n],
  \]
  where $R_\epsilon = \left\{ r \in R : \lvert r \rvert \leq \varepsilon \right\}$ and $\res{R}_\varepsilon
  = \pb{R} / R_\varepsilon$, with kernel
  \[
    \ker \tau_\varepsilon = \left\{ f \in T_n : \lvert f \rvert \leq \varepsilon \right\}.
  \]
  $\tau_\varepsilon (g)$ is a unitary polynomial in $x_n$ of degree $s$, and therefore we can apply
  Euclidean division (in $(\res{R}_\varepsilon [x_1, \cdots, x_{n-1}]) [x_n]$) on it.
  %% I need to work out why this is the case...?
  This means for all $f \in \pb{T}_n$ we can find $q \in \pb{T}_n$ and $r \in \pb{T}_{n-1}[x]$ with
  degree $ < s$ such that
  \[
    \tau_\varepsilon (f) = \tau_\varepsilon (q) \tau_\varepsilon (g) + \tau_\varepsilon (r)
  \]
  or equivalently,
  \[
    \lvert f - (qg + r) \rvert \leq \varepsilon
  \]
  since this will be in the kernel of $\tau (\varepsilon)$.
  Hence, for all $f \in T_n$, there exists $b \in B$ such that
  \[
    \lvert f - b \rvert \leq \varepsilon \lvert f \rvert,
  \]
  meaning $B$ is $\varepsilon$ dense in $T_n$ and therefore dense.
  This completes the proof, as $B$ closed and dense implies $B = T_n$.
\end{proof}

\section{Weierstrass preparation}

Earlier we claimed that Weierstrass preparation follows from Weierstrass division. We now show that
this is the case.

\begin{lemma}
  A series $f \in T_n$ with $\lvert f \rvert = 1$ is a unit in $T_n$ if and only if $\lvert f (0) \rvert = 1$
  and $\lvert f - f(0) \rvert < 1$. Thus, $f$ is a unit if and only if $\tilde{f}$ is a unit;
  i.e. a constant since $\res{T}_{n} \isom \tilde{R}[x_1, \cdots, x_{n}]$.
\end{lemma}

\begin{proof}
  Since $\lvert \cdot \rvert$ is a valuation, $f$ is a unit if and only if it is a unit in $\pb{T}$.
  This is equivalent to $a_{(0, \cdots, 0)}$ being a unit in $\pb{R}$ and $a_{(i_1,\cdots, i_n)} \in \tn{R}$
  for all other indices. This is the assertion.
\end{proof}

%% This depends on Proposition 1.4.2 3; trying to work out the proofs for this right now.

\begin{theorem}
  Let $g \in T_n$ be $x_n$-distinguished of degree $s$, for some natural number $s$. Then there is a
  unique monic polynomial $w \in T_{n-1}[x_n]$ of degree $s$ and a unique unit $e \in T_n$ such that
  \[
    g = e w
  \]
  with $\lvert w \rvert = 1$, so that $w$ is $x_n$-distinguished of degree $s$.
\end{theorem}

\begin{proof}
  By Weierstrass division, there exists unique $e' \in T_n$ and $r' \in T_{n-1} [x_n]$ with degree $r' < s$
  such that
  \[
    x_n^s = e' g + r.
  \]
  Then, defining
  \[
    w = x^s_n - r
  \]
  we see $w$ is monic of degree $s$ and $w = e' g$.

  It thus suffices to show that $e'$ is a unit, because then $g = w e^{-1}$ with the wanted properties.

  Using the estimates in Weierstrass division we have
  \[
    \lvert r \rvert \leq \lvert x_n^s \rvert = 1,
  \]
  and so we see that $\lvert w \rvert = 1$ and $x_n$ is distinguished of degree $s$.
  Now assume $\lvert g \rvert = 1$. From $\tilde{w} = \tilde{e'} \tilde{g}$, we conclude that $\tilde{e'}$
  is a unit in $\res{T}_{n-1}$, and since $\res{T}_{n-1} \isom \tilde{R}[x_1, \cdots, x_{n-1}]$ this means
  $\tilde{e'}$ is a constant.
  A fortiori this means it is a unit in $\res{T_n}$ and so invoking the above lemma it is a unit in $T_n$
  and we are done with the existence part.

  For uniqueness, suppose $g = e w$, define $r = x^s_n - w$. Then
  \[
    x_n^s = e^{-1} g + r
  \]
  is a valid representation satisfying the conditions for Weierstrass division and therefore $e$ and
  $w$ are unique, and moreover so is $w$.
\end{proof}


\section{Aside lemmas to fit in}

\begin{lemma}
  Let $A$ be a complete topological ring, then each element of the form $e = 1 - y$, where $y$ is topologically
  nilpotent, is a unit in $A$.
  We have $ e{-1} = \sum_{n=0}^\infty = 1 + z$, where $z$ is topologically nilpotent.
\end{lemma}

\begin{proof}
  The ring being complete we can set
  \[
    v = 1 + z, \quad \text{where} \quad z = \sum_{n=1}^\infty y^n.
  \]
  Since topological nilpotents are closed and $y^n$ is topologically nilpotent for all $n \geq 1$, we
  get $z$ is topologically nilpotent. We have
  \[
    e v = (1 - y) \sum_{n = 0}^{\infty} y^n = 1
  \]
  i.e. $v = e^{-1}$.
\end{proof}


\begin{lemma}
  Let $A$ be a complete topological space. Then a powerbounded element $a$ is a unit if and only if
  its residue class $\tilde{a} \in \tilde{A}$ is a unit.
\end{lemma}

\begin{proof}
  If $a$ is a uni then it is immediate that $\tilde{a}$ is a unit (take its inverse and reduce to
  its residue class, then they are inverses in $\tilde{A}$).

  To show the converse, consider a power bounded $b$ with $\tilde{a} \tilde{b} = \tilde{1}$. This
  means $a b = 1 - x$ for some topologically nilpotent $x$. By the previous lemma we conclude that
  $ab$ is a unit and hence $a$ is a unit.
\end{proof}

\begin{lemma}
  If $R$ is a complete normed space and $c$ is a positive real number. Then $R_c\lvert x \rvert$
  (restricted power series for parameter $c$) is complete with respect to the gauss norm.
\end{lemma}

\begin{proof}
  Let $(f_i) = (s\sum_{n = 0}^\infty a_{n,i} x^n)_i$ be a Cauchy sequence in $R\lvert x \rvert$.
  Since
  \[
  \lvert a_{n,i+1} - a_{n,i} \rvert * c^n \leq \lvert f_{i+1} - f_i \rvert
  \]
  (by writing out the definition of gauss norm of $f_{i+1} - f_i$) for fixed $n$, we see that each
  sequence $(a_{n,i})_i$ is a Cauchy sequence in $R$.

  Let $a_n$ be a limit of this sequence and set
  \[
    f = \sum_{n = 0}^\infty a_n x^n.
  \]
  We want to show that $f$ is a restricted power series for parameter $c$ and that
  \[
    \lim_{i \to \infty} \lvert f - f_i \rvert = 0.
  \]
  We may assume $\lvert f_j - f_i \rvert \leq 1 / i$ for all $j \geq i$.
  From all $\lvert a_{n,j} - a_{n,i} \rvert * c^n \leq 1/i$ we deduce by continuity of the norm on $R$
  that $\lvert a_n - a_{n,i} \rvert * c^n \leq 1 / i$ for all $n \geq 0$ and all $i$. For $n$ big enough
  we have $\lvert a_{n,i} \rvert *c^n < 1 /i$ since $f_i$ is restricted.
  Therefore, $\lvert a_n \rvert * c^n \leq 1/i$ and we get $f$ is restricted.

  Furthermore, we have
  \[
    \lvert f - f_i \rvert = \max \lvert a_n - a_{n,i} \rvert c^n \leq 1/i
  \]
  and, so we see $f$ is the limit of $f_i$.
\end{proof}

\begin{lemma}
  If $R$ is complete so is $\pb{R}$ the ring of power bounded elements.
\end{lemma}

\begin{proof}
  $\pb{R}$ is closed in a complete space.
\end{proof}

\begin{lemma}
  We have $\pb{R} \langle x \rangle = \pb{R \langle x \rangle}$ and
  $\tn{R} \langle x \rangle = \tn{R \langle x \rangle}.$
\end{lemma}

\begin{proof}
  %% This needs to be copied out... concern is that this may need to be an isometry...
  %% which is not hard but just more work
  %% perhaps I can see if the proof really needs it
  %% Otherwise perhaps I just want this as a statement saying that something in the LHS
  %% is powerbounded or topologically nilpotent respectively in the whole space
  %% which seems like it would be better
  %% I can then also have this as an iff if I want
\end{proof}


\begin{lemma}
  Let $R$ be complete. An element $f = \sum_{n = 0}^\infty a_n x^n \in \pb{R} \langle x \rangle$ (that
  is every coefficient is power bounded) is a unit in $\pb{R} \langle x \rangle$ if and only if $a_0$
  is a unit and $a_n$ is topologically nilpotent for all $n > 0$.
\end{lemma}

\begin{proof}
  Since $\pb{R}$ is complete we have $\pb{A} \langle x \rangle$ is also complete. Then by an earlier
  lemma we have that $f$ is a unit in $\pb{R} \langle x \rangle$ if and only if its residue class
  $ \tilde{f} = \sum_{n=0}^\infty \tilde{a}_n x^n$ is a unit.
  Moreover, we have that $a_0$ is a unit in $\pb{R}$ if and only if $\tilde{a_0}$ is a unit.

  Therefore, it is enough to show that any unit in the residue class is a unit in $\tilde{R}$, that
  is its a constant power series.
  %% This statement may have some nuance...
  %% For example, we want the reduction \tilde{R \langle x \rangle } to be \tilde{R} \langle x \rangle
  %% This may be done seperately or see if it can be avoided.
\end{proof}
